\documentclass[11pt]{article}
\usepackage[margin=0.88in]{geometry}
\usepackage{amsmath,amssymb}
\usepackage{microtype}
\usepackage{xcolor}
\usepackage{hyperref}
\hypersetup{colorlinks=true,urlcolor=blue!55!black,citecolor=blue!55!black}
\newcommand{\Hs}{\mathcal H_s}
\newcommand{\Hl}{\mathcal H_\ell}
\newcommand{\Mult}{\operatorname{Mult}}
\title{The mixed complete-Pick interpolation question\\was answered negatively by arXiv:2210.14477}
\author{Literature-status packet for arXiv:1701.04885}
\date{11 August 2026}

\begin{document}
\maketitle

\begin{center}
\fbox{\parbox{0.92\linewidth}{\textbf{Status: literature already answered.}
Tsikalas explicitly restates the Aleman--Hartz--McCarthy--Richter question,
constructs counterexamples to its proposed equivalence, and proves the
correct replacement criterion.  This is a retrieval result, not a new
counterexample produced by the present run.}}
\end{center}

\section*{Original question}

Aleman, Hartz, McCarthy, and Richter ask in the final paragraph of
Section~6, ``Applications and Questions,'' page 21 of the arXiv source
PDF~\cite{AHMR}:

\begin{quote}
If $s$ is a normalized complete Pick kernel and $\ell=gs$ for some
positive-semidefinite kernel $g$, is a sequence $(\lambda_i)\subseteq X$
interpolating for $\Mult(\Hs,\Hl)$ if and only if it is a Carleson sequence
for $\Hs$ and is $\Hl$-separated?
\end{quote}

The source had proved this under additional hypotheses, notably when the
target kernel is a power of a complete Pick kernel, but left the general
statement open.

\section*{Decisive later answer}

Georgios Tsikalas, \emph{Interpolating sequences for pairs of spaces},
arXiv:2210.14477; J. Funct. Anal. 285 (2023), 110059~\cite{Tsikalas}, explicitly
labels the same statement as Question~1.1 and credits it to the four source
authors.  The abstract likewise says that the paper answers their question.

The answer is \textbf{no in general}.  Supporting Theorem~1.5 constructs a
kernel $\ell$ with a normalized complete Pick factor $s$ such that, for every
$n\ge2$, there is a sequence satisfying the $\Hs$ Carleson condition and
$n$-weak separation by $\ell$, but not $(n+1)$-weak separation and hence not
$\Mult(\Hs,\Hl)$ interpolation.  Taking $n=2$ gives exactly a sequence that is
pairwise (ordinary weakly) separated and $\Hs$-Carleson but is not
interpolating, disproving the proposed equivalence.

Supporting Theorem~1.4 gives the corrected full characterization:
\[
(\lambda_i)\text{ is }\Mult(\Hs,\Hl)\text{-interpolating}
\quad\Longleftrightarrow\quad
\begin{cases}
(\lambda_i)\text{ is }\Hs\text{-Carleson},\\
(\lambda_i)\text{ is }n\text{-weakly separated by }\ell
\text{ for every }n\ge2.
\end{cases}
\]
Here $n$-weak separation uniformly bounds the distance of any one normalized
$\ell$-kernel from the span of any other $n-1$ kernels.  Ordinary separation
is only the case $n=2$.

\section*{Provenance and scope}

This is an unambiguous \texttt{literature\_already\_answered} match: the later
author explicitly cites and restates the original question, gives a negative
counterexample theorem, and replaces the false criterion by an exact one.
The later paper also identifies important classes where ordinary separation
does suffice, via its automatic-separation property.  Those positive subclasses
do not change the negative answer to the source's unrestricted question.

The packet contains the original paper as \texttt{source\_paper.pdf} and the
decisive later paper as \texttt{supporting\_paper\_2210.14477.pdf}.  A matching
run-ledger record accompanies this status packet.

\begin{thebibliography}{9}
\bibitem{AHMR}
A. Aleman, M. Hartz, J. E. McCarthy, and S. Richter,
\emph{Interpolating sequences in spaces with the complete Pick property},
Int. Math. Res. Not. IMRN 2019, no. 12, 3832--3854;
\href{https://arxiv.org/abs/1701.04885}{arXiv:1701.04885}.

\bibitem{Tsikalas}
G. Tsikalas,
\emph{Interpolating sequences for pairs of spaces},
J. Funct. Anal. 285 (2023), 110059;
\href{https://doi.org/10.1016/j.jfa.2023.110059}{doi:10.1016/j.jfa.2023.110059};
\href{https://arxiv.org/abs/2210.14477}{arXiv:2210.14477}.
\end{thebibliography}

\end{document}
